\chapter{埋め込みとポジショナルエンコーディング}
\label{ch:embedding}

この章では、トークンIDを密集ベクトルに変換する\keyterm{埋め込み（embedding)}と、シーケンスの順序情報をモデルに注入する\keyterm{ポジショナルエンコーディング（positional encoding)}を実装します。

\section{なぜ埋め込みが必要なのか}

第3章では、トークナイザーが「Hello」を整数ID 247に変更しました。しかし、この定数をそのままニューラルネットワークに入れると問題が生じる。まず、整数247と248が「同様の単語」であるという保証はありません。 「cat」は247、「dog」は248、または「zebra」は248です。第二に、整数を乗算または加算する演算は意味を持たない。 「cat + 1 = dog」は成立しません。

解決策は、各トークンIDを\keyterm{密集ベクトル}に変換することです。語彙サイズが$V$で埋め込み次元が$d$の場合、埋め込みはサイズ$V \times d$の学習可能な行列$E$です。トークンID$i$が入ると、$E$の$i$行目が取得されます。この行ベクトルがそのトークンの「意味のある表現」になります。

\begin{infobox}[なぜ「密集」なのか？]
ワンホットエンコーディングは、語彙サイズ$V$のスパースベクトルでトークンを表現します。たとえば、10000語彙の場合、ID 247は246個のゼロと1個の1、残りの0からなるベクトルです。これを$E$と乗算すると、最終的に$E$の247番目の行を取得するのと同じです。埋め込みは、この「ワンホット$\times$行列」をループ操作として実行する効率的な実装です。

学習が進むと、同様の意味の単語がベクトル空間に近づきます。 「cat」と「dog」は近く、「cat」と「car」は遠いです。これが埋め込みの魔法である: 学習前はランダムベクトルであったが、学習後には意味を込めた座標となる。
\end{infobox}

\section{埋め込みの数学的定義}

埋め込み行列$E \in \mathbb{R}^{V \times d}$がある場合、トークンID$i$の埋め込みは$E[i] \in \mathbb{R}^d$です。これは単に$E$の$i$番目の行を選択する演算であり、微分可能である (backward 時選択された行だけグラデーションを受け取る)。

バッチ入力$\mathbf{x} \in \mathbb{Z}^{B \times L}$（バッチ$B$、シーケンス長$L$）の場合、埋め込み出力は$\mathbf{X} \in \mathbb{R}^{B \times L \times d}$になります。

\section{Token Embedding 実装}

PyTorchの\code{nn.Embedding}は正確にこのlookup操作を実行します。私たちはそれを薄くラップして初期化ロジックを追加します。

\begin{lstlisting}[language=Python]
import math
import torch.nn as nn

class TokenEmbedding(nn.Module):
    def __init__(self, vocab_size: int, d_model: int, pad_id: int = 0):
        super().__init__()
        self.vocab_size = vocab_size
        self.d_model = d_model
        self.pad_id = pad_id
        self.embedding = nn.Embedding(vocab_size, d_model, padding_idx=pad_id)
        self._init_weights()

    def _init_weights(self):
        # 正規分布 N(0, 1/sqrt(d_model)) 初期化(GPT-2スタイル)
        nn.init.normal_(self.embedding.weight, mean=0.0,
                        std=1.0 / math.sqrt(self.d_model))
        # パディングトークンは0
        with torch.no_grad():
            self.embedding.weight[self.pad_id].fill_(0.0)

    def forward(self, token_ids: torch.Tensor) -> torch.Tensor:
        # [batch, seq] -> [batch, seq, d_model]
        return self.embedding(token_ids)
\end{lstlisting}

\begin{notebox}[初期化の重要性]
ディープラーニングにおける重みの初期化は学習成否に左右される。大きすぎるとアクティブ値が爆発し、小さすぎるとグラデーションが消えます。 GPT-2は標準偏差$1/\sqrt{d\_{model}}$の正規分布を使用し、これは初期埋め込みベクトルのnormが約1になるような合理的な選択です。

なぜあなたは$1/\sqrt{d}$ですか？$d$次元ベクトルの各要素が$\mathcal{N}(0, \sigma^2)$の場合、ベクトルnormのコメント値は$\sqrt{d} \cdot \sigma$です。$\sigma = 1/\sqrt{d}$なら norm が 1 近くになる。これにより、埋め込みベクトルのスケールがレイヤ出力に似ており、学習が安定します。
\end{notebox}

\subsection{padding\_idxの役割}

\code{nn.Embedding}の\code{padding\_idx}パラメータは、パディングトークンの埋め込みを常にゼロに保ちます。これにより：
\begin{itemize}
\item パディング位置の出力が0になり、アテンション計算に影響を与えません。
\item backward 時、パディング行のグラデーションが 0 になり、重みが更新されない。
\end{itemize}

\section{ポジショナルエンコーディング}

トランスはアテンションだけでシーケンスを処理しますが、アテンション自体には「順序」という概念はありません。入力を混ぜても出力は（再配列されたまま）同じ。これを\keyterm{順序不変性（permutation invariance)}という。自然言語では、順序が意味を決定するので（「犬が人を尋ねた」対「人が犬を尋ねた」）、順序情報を明示的に注入する必要があります。

\subsection{直観: なぜ順序が重要なのか?}

アテンションは、「各トークンが他のトークンにどれだけ注意を払うか」を学びます。しかし、「最初のトークン」と「最後のトークン」の区別はありません。文章「私はご飯を食べました」から「食べた」が与えられ、「私は」とつながるべきですが、アテンションは「5番目のトークンが1番目のトークンを見る」という情報がなければこの接続を学習することは困難です。

ポジショナルエンコーディングは、各位置に固有の「位置ベクトル」を追加し、モデルが位置を区別できるようにします。

\subsection{Sinusoidal 位置エンコーディング（Vaswani et al. 2017)}

オリジナルトランスフォーマー論文が提案した方式。サインとコサイン関数で位置をエンコードします。

\begin{formula}[Sinusoidal Positional Encoding]
\[
PE\_{(pos, 2i)} = \sin\left(\frac{pos}{10000^{2i/d\_{model}}}\right)
\]
\[
PE\_{(pos, 2i+1)} = \cos\left(\frac{pos}{10000^{2i/d\_{model}}}\right)
\]
ここで、$pos$は位置、$i$はディメンションインデックスです。偶数次元はsin、奇数次元はcos。
\end{formula}

この方式の利点は、学習可能なパラメータがないことです。また、相対位置に対する線形関係が成立し、モデルが容易に学習できる。$\sin(\alpha + \beta) = \sin\alpha\cos\beta + \cos\alpha\sin\beta$なので、位置$pos + k$の符号化は、位置$pos$の符号化の線形変換として表される。

\begin{lstlisting}[language=Python]
class PositionalEmbedding(nn.Module):
    def __init__(self, d_model: int, max_len: int = 512, mode: str = "learned"):
        super().__init__()
        self.mode = mode
        if mode == "sinusoidal":
            pe = torch.zeros(max_len, d_model)
            position = torch.arange(0, max_len).unsqueeze(1).float()
            div_term = torch.exp(
                torch.arange(0, d_model, 2).float() * (-math.log(10000.0) / d_model)
            )
            pe[:, 0::2] = torch.sin(position * div_term)
            pe[:, 1::2] = torch.cos(position * div_term)
            self.register_buffer("pe", pe.unsqueeze(0))  # [1, max_len, d_model]
        else:  # learned
            self.pe = nn.Embedding(max_len, d_model)
            nn.init.normal_(self.pe.weight, mean=0.0, std=1.0/math.sqrt(d_model))

    def forward(self, seq_len: int) -> torch.Tensor:
        # [1, seq_len, d_model] 返す
        if seq_len > self.max_len:
            raise ValueError(f"Sequence length {seq_len} exceeds max_len {self.max_len}")
        if self.mode == "sinusoidal":
            return self.pe[:, :seq_len, :]
        else:
            positions = torch.arange(seq_len, device=self.pe.weight.device)
            return self.pe(positions).unsqueeze(0)
\end{lstlisting}

\begin{figure}[htbp]
\centering
\begin{tikzpicture}[
  scale=0.7,
  every node/.style={font=\small\sffamily}
]
  % sin波形
  \draw[inkblue, line width=0.5pt, ->] (0, 0) -- (8, 0) node[right] {pos};
  \draw[inkblue, line width=0.5pt, ->] (0, -1.2) -- (0, 1.2) node[above] {value};
  \draw[inkred, line width=1pt, domain=0:7, samples=50] plot (\x, {sin(\x*50)});
  \node[inkred, font=\scriptsize] at (7, 0.8) {dim 0 (sin)};
  \draw[inkblue, line width=1pt, domain=0:7, samples=50] plot (\x, {cos(\x*50)});
  \node[inkblue, font=\scriptsize] at (7, -0.8) {dim 1 (cos)};
\end{tikzpicture}
\caption{Sinusoidal 位置エンコーディングの最初の2次元. sinと cosこの位置によって周期的に変化する.}
\label{fig:sinusoidal}
\end{figure}

\subsection{学習可能な位置埋め込み（GPT-2スタイル)}

ＧＰＴ－２、ＧＰＴ－３などは、位置ごとに学習可能な埋め込みベクトルを使用する。このアプローチは、十分なデータがある場合はより柔軟ですが、学習時に見たことのない長さ（例：学習時1024でしたが、推論時2048）には対応することは困難です。

\begin{notebox}[Sinusoidal vs Learned]
\begin{itemize}
\item\textbf{Sinusoidal}：パラメータなし、extrapolation可能（学習時の本長より長いシーケンスにも対応）。しかし、実際のパフォーマンスは学習可能な方法より少し低下する傾向。
\item\textbf{Learned}：パフォーマンスは少し優れていますが、固定長。 extrapolationの難しさ。
\end{itemize}
最新のモデルは、RoPE（Rotary Position Embedding）を使用して2つの方法の利点を組み合わせています。 2巻5章で扱う。
\end{notebox}

\section{埋め込み合成}

トークンの埋め込みと位置の埋め込みを加えて、最終的な入力表現を作成します。 Vaswani et al。は、埋め込みに$\sqrt{d\_{model}}$を掛けてスケールを調整することを提案しました。これは、位置埋め込みとスケールを同様に維持するためです。

\begin{lstlisting}[language=Python]
class EmbeddingStack(nn.Module):
    def __init__(self, vocab_size, d_model, max_len=512,
                 pad_id=0, pos_mode="learned", dropout=0.1):
        super().__init__()
        self.token_emb = TokenEmbedding(vocab_size, d_model, pad_id=pad_id)
        self.pos_emb = PositionalEmbedding(d_model, max_len, mode=pos_mode)
        self.dropout = nn.Dropout(dropout)
        self.scale = math.sqrt(d_model)

    def forward(self, token_ids: torch.Tensor) -> torch.Tensor:
        batch, seq = token_ids.shape
        tok = self.token_emb(token_ids) * self.scale  # スケーリング
        pos = self.pos_emb(seq)
        return self.dropout(tok + pos)
\end{lstlisting}

\begin{figure}[htbp]
\centering
\begin{tikzpicture}[
  every node/.style={font=\small\sffamily},
  box/.style={draw=inkblue, fill=inkblue!8, rounded corners=2pt, minimum width=2.5cm, minimum height=0.7cm},
  sum/.style={draw=inkred, fill=inkred!10, circle, minimum size=0.6cm}
]
  \node[box] (ids) at (0, 0) {token IDs};
  \node[box, below=0.6cm of ids] (tok) {Token Embedding};
  \node[box, right=2cm of ids] (pos) {Position};
  \node[box, below=0.6cm of pos] (posemb) {Positional Embedding};
  \node[sum, below=0.8cm of $(tok)!0.5!(posemb)$] (add) {$+$};
  \node[box, below=0.6cm of add] (out) {Dropout $\to$ output};
  \draw[-{Stealth[length=4pt]}, inkblue] (ids) -- (tok);
  \draw[-{Stealth[length=4pt]}, inkblue] (pos) -- (posemb);
  \draw[-{Stealth[length=4pt]}, inkblue] (tok) -- (add);
  \draw[-{Stealth[length=4pt]}, inkblue] (posemb) -- (add);
  \draw[-{Stealth[length=4pt]}, inkblue] (add) -- (out);
\end{tikzpicture}
\caption{EmbeddingStack 構造. トークン埋め込みと位置埋め込みを追加した後にドロップアウトを適用する.}
\label{fig:embedding-stack}
\end{figure}

\subsection{なぜスケーリングが必要なのか?}

トークン埋め込みは$\mathcal{N}(0, 1/d)$に初期化されており、元素の分散が$1/d$です。位置埋め込みの元素は$[-1, 1]$の範囲で分散が約$1/2$である。スケーリングなしで加算すると、位置情報はトークン情報を圧倒します。

$\sqrt{d}$を掛けると、トークン埋め込みの分散が$1$になり、位置埋め込みとバランスが取れます。

\section{Layer Normalization はじめに}

トランスで学習を安定させるために\keyterm{レイヤー正規化（Layer Normalization)}を使用してください。次章のアテンションでも登場するのでここで紹介する。

\begin{formula}[Layer Normalization]
\[
\text{LN}(x) = \gamma \cdot \frac{x - \mu}{\sqrt{\sigma^2 + \epsilon}} + \beta
\]
ここで、$\mu, \sigma^2$は最後の次元に沿って計算された平均と分散です。$\gamma, \beta$は学習可能なパラメータです。
\end{formula}

PyTorchでは\code{nn.LayerNorm(d\_model)}を1行で使用します。 2巻ではLayerNormの変形であるRMSNormも扱う。

\begin{infobox}[LayerNorm vs BatchNorm]
BatchNormは配置次元に沿って正規化し、学習と推論の際に統計が異なる問題がある。 LayerNormは各サンプル内で正規化し、バッチサイズに関係なく動作します。これがLLMでLayerNormを書く理由だ。 バッチサイズ1の推論でも安定して動作する。
\end{infobox}

\section{テスト}

埋め込みモジュールのコアテストは、出力形状の確認と勾配の流れの確認です。

\begin{lstlisting}[language=Python]
def test_token_embedding_shape():
    emb = TokenEmbedding(vocab_size=100, d_model=32)
    ids = torch.randint(0, 100, (2, 8))  # [batch=2, seq=8]
    out = emb(ids)
    assert out.shape == (2, 8, 32)

def test_positional_sinusoidal():
    pe = PositionalEmbedding(d_model=32, max_len=64, mode="sinusoidal")
    out = pe(seq_len=16)
    assert out.shape == (1, 16, 32)

def test_embedding_stack():
    stack = EmbeddingStack(vocab_size=100, d_model=32, max_len=64, dropout=0.0)
    ids = torch.randint(0, 100, (2, 8))
    out = stack(ids)
    assert out.shape == (2, 8, 32)

def test_positional_exceeds_max_len_raises():
    pe = PositionalEmbedding(d_model=32, max_len=32, mode="learned")
    with pytest.raises(ValueError):
        pe(seq_len=64)  # max_len 超過
\end{lstlisting}

\section{練習: 埋め込み可視化}

\code{scripts/tiny\_train.py}で学習された埋め込みをPCAで2次元に投影すると、単語がどのように分布するかがわかります。十分な学習の後、似たような単語が近くに集まることを確認できます。

\section{要約}

\begin{itemize}
\item 埋め込みは、トークンIDを密集ベクトルに変換する学習可能なlookupテーブルです。
\item 初期化は$N(0, 1/\sqrt{d\_{model}})$を使用し、パディングトークンは0に設定します。
\item ポジショナルエンコーディングはトランスに順序情報を注入します。 sinusoidalとlearnedの2つの方法があります。
\item トークン埋め込みと位置埋め込みを追加し、ドロップアウトを適用して最終入力表現を作成します。
\item $\sqrt{d\_{model}}$スケーリングで2つの埋め込みの分散をバランスさせます。
\item LayerNormはバッチサイズに依存しない正規化で、LLM学習の安定化に不可欠です。
\end{itemize}

\section{練習問題}

\begin{enumerate}
\item 埋め込みディメンション$d$が小さすぎる場合（例：8）、どのような問題が発生しますか？大きすぎる場合（例：4096）、どのような問題が発生しますか？
\item Sinusoidal位置エンコーディングで\code{10000}という定数を100に変更するとどうなりますか？
\item トークン埋め込みに$\sqrt{d\_{model}}$を掛けないと、学習はどのように異なりますか？
\item 学習可能な位置埋め込みの\code{max\_len}を1024に設定して学習した後、推論時に2048の長さのシーケンスを入れるとどうなりますか？
\item LayerNormの$\epsilon$値を$10^{-5}$から$10^{-1}$に変更すると、出力がどのように変わりますか？
\end{enumerate}

次の章では、この埋め込みを入力として受け取り、「各トークンが他のトークンにどれだけ注意を払うか」を学習するアテンションメカニズムを実装します。
